\documentclass{article}
\usepackage[utf8]{inputenc}
\usepackage[T1]{fontenc}
\usepackage[margin=2cm]{geometry}
\usepackage{hyperref}
\sloppy
\begin{document}
\section*{Historical source}
Encyclopaedia Britannica scan, source row 1649. Parse 3.5 extracted the text. Super 3.5 VL corrected it from an enlarged crop of each block. It can still contain recognition errors.
Page headers, footers, and page numbers are removed.

\subsection*{Untitled continuation (E0142)}
\textit{Continues an entry from an earlier page.}

Chivalry, and they extend their influence equally through the whole system. Difference of armour, and different modes of military discipline, produce distinction of orders among those who practise the arts of war; while other circumstances, originating from the same general causes, occasion similar changes to take place amid the scenes of peace.

6. The distinction introduced into the military order by the use of cavalry.

None of the new distinctions which are introduced among men, with respect to the discipline and conduct of war, in consequence of the acquisition of property and the invention of arts, is more remarkable than that occasioned by the use of horses in military expeditions, and the training of them to the evolutions of the military art. Fire-arms, it is true, give to those who are acquainted with them a greater superiority over those to whom their use is unknown than what the horseman possesses over him who fights on foot. But the use of fire-arms is of such importance in war, and the expense attending it so inconsiderable, that wherever they have been introduced, they have seldom been confined to one particular order in an army; and therefore they produce indeed a remarkable, though transient, distinction among different nations; but establish no permanent distinctions in the armies of any one nation. But to maintain a horse, to equip him with costly furniture, to manage him with dexterity and vigour, are circumstances which have invariably produced a striking and conspicuous distinction among the military order, wherever bodies of cavalry have been formed. The Roman equites, who, though they became at length a body of usurers and farmers-general, were originally the only body of cavalry employed by the state, occupied a respectable rank between the senators and the plebeians; and the elegance and humanity of their manners were suitable to their rank. In ancient Greece, and in the celebrated monarchies of Asia, the same distinction prevailed at a similar period.

Military dis-
tinctions among the
ancient Germans.

Another fact worthy of notice reflecting the manners of the barbarians of Germany before they clasped arms among the savage warriors of Scythia and Germany, as well as among the Greeks or Romans. From the valuable treatise of Tacitus de moribus Germanorum, we learn, that among the German warriors a distinction somewhat of this nature did actually subsist; not so much indeed a distinction between the warrior who fought on horseback and those who fought on foot, as between those whom vigour of body and energy of mind enabled to brave all the dangers of war, and such as, from the imbecility of youth, the infirmities of age, or the natural inferiority of their mental and bodily powers, were unequal to scenes of hardship and deeds of valour. The youth was not permitted to take arms and join his warlike countrymen in their military expeditions whenever he himself thought proper. There was a certain age before which he could not be invested with armour. When he had attained that period, if not found deficient in strength, activity, or courage, he was formally honoured with the field and the lance, called to the duties, and admitted to all the privileges, of a warrior.

billified themselves in the cultivated provinces of the Chirary, Roman empire is, that their women, contrary to what we and among many other rule nations, were treated with a high degree of refuted. They did not guess 8 rally rise via the men in deeds of values but they share of the animalist them by their exhordiations to diffusion and warn as freedom in the field, and virtue of expertise were evenge the confided with a faced veneration, as endowed with propried powers, capable to forced eternal aid in the form of fluitty, and even to influence the all of the dieties. Hence, though domestic duties were their peculiar province, yet they were not hardly tested the confined to a state of slavery. There appears in died a striking analogy between the condition of the women among the endold dies of Spain and the rank which, they occupied among the warkite cantons of Germany. Perhaps indyed, the Germans were fill more honorable than the Spartan woman, as they were taught to weld the magic weapons of imperfection, which life Greece were propagated to the prefix.

It appears, therefore, that, in the forests of Germany at least, if not in the more northern regions of Asia and Europe, the conquerors of the Roman empire, before they penetrated into its provinces, treated their women with a degree of respect unknown to most of the nations of antiquity; that the character of the warrior was likewise highly honourable, being understood to unite all those qualities which were in the highest estimation; and that it was only at a particular age, and with certain forms, that the youth were admitted to bear arms.

When those nations sailed from their defects and forfeits, over-ran the Roman empire, and established themselves in its provinces, the change which took place on their circumstances was remarkable; and by a natural influence, it could not but produce an equally remarkable change on their habits, customs, and manners. The great outlines might fill remain; but they could not now fail to be filled up in a different manner. Here, however, the records of history are

pecially imperfect. We have no Caesar or Tacitus to supply facts or direct our reasonings; the Gothic nations had not yet learned to read and write; and the Romans were so depraved under the sense of their own miseries, as to be negligent of the changes which happened around them. But as soon as the light of history begins again to dawn, we find that the leading features of the barbarian character were not effaced, but only modified in a particular manner, in consequence of their mixing among a more polished people, becoming acquainted with the luxuries of life, and acquiring extensive power and property.

Those who fought on horseback now began to be distinguished with peculiar honours. The manners of the warrior too were become more cultivated, and his spirit more humane. Leisure and opulence, with the influence of a polished people, even though in a state of slavery, taught those barbarians to aspire after more refined pleasures and more splendid amusements than those which they had been before satisfied with. The influence of Christianity too, which, though greatly corrupted, was still favourable to the social happiness of mankind, concurred to polish their manners and exalt their character. Hence, in the end of the tenth and in the beginning of the eleventh century, we fee


\section*{Source record}
\begin{description}
\item[Scan ID] \texttt{britannica:309f765b1ecc83e0a7f7762e44b4a927a7555443:train:1649}
\item[Image file] \texttt{page-1649-e6a1fc569876.jpg}
\item[Image SHA-256] \texttt{e6a1fc5698767dd40265a849aa559e35ef251daca6cf3e82cf486dd083a4e3e0}
\item[Dataset] NationalLibraryOfScotland/encyclopaedia\_britannica\_illustrated at revision \texttt{309f765b1ecc83e0a7f7762e44b4a927a7555443} (CC0-1.0)
\item[Extraction] model \texttt{nvidia/nemotron-3.5-parse}, raw reply \texttt{extract/raw/page-1649-e6a1fc569876.json}
\item[Raw reply SHA-256] \texttt{72fb98bbf2b620f07ba563b39e097d119624c96a18a558a5e57d526bffcbb98c}
\item[Correction] model \texttt{nvidia/Nemotron-3.5-Super-VL}, block crops, raw replies \texttt{correct/raw/crop\_nothink/page-1649-e6a1fc569876-b*.json}; parse\_fallback: 1, super\_vl: 9
\item[Running header] CH I 671 CHI
\end{description}
\end{document}
